\section*{Chapitre \ref{ch:b1:radiation-heat-transport} --- Le rayonnement et le transport de chaleur}
% values: calc/ch27.py

\begin{solution}{exo:b1:radiation-heat-transport:1}
(a)~$S_0/\pi = 1361/\pi = \qty{433.2}{W.m^{-2}}$. (b)~À l'équateur à un équinoxe
$\bar Q = (S_0/\pi)(a/r)^2$, donc $(a/r)^2 = 436.7/433.2 = 1.0080$ et $r/a =
1.0080^{-1/2} = 0.9960$~: la Terre est \qty{0.4}{\percent} plus près du Soleil que
sa distance moyenne. (c)~$\bar Q = 433.2 \times 0.9683 \times \cos\ang{23.44} =
\qty{384.9}{W.m^{-2}}$.
\end{solution}

\begin{solution}{exo:b1:radiation-heat-transport:2}
(a)~$415.6/172.4 = 2.41$. (b)~La zone entre les parallèles $\pm\varphi$ a pour
aire $2\pi a_\oplus^2 \times 2\sin\varphi$, une fraction $\sin\varphi$ de la sphère~:
$\sin\ang{37.0} = 0.602$, \qty{60}{\percent}. (c)~Sur un axe en \omterm{def:b1:spherical-astronomy:geographic}{latitude}, les
bandes près des pôles paraissent aussi larges que celles près de l'équateur, mais elles couvrent bien
moins d'aire~: la zone au-delà de \ang{60} ne fait que $1 - \sin\ang{60} =
\qty{13.4}{\percent}$ de la surface. En fonction de $\sin\varphi$, les tropiques remplissent la plus grande part
de l'axe.
\end{solution}

\begin{solution}{exo:b1:radiation-heat-transport:3}
(a)~$T_e = (\text{OLR}/\sigma)^{1/4}$~: \qty{256.4}{K}, \qty{241.3}{K} et
\qty{254.5}{K}. (b)~Lumière solaire absorbée~: $287.7/75.8 = 3.80$~; OLR~: $244.9/192.3 =
1.27$. (c)~Le \omterm{def:b1:radiation-heat-transport:net-radiation}{rayonnement net} suit la lumière solaire absorbée~: positif et grand dans les
tropiques, fortement négatif aux hautes \omterm{def:b1:spherical-astronomy:geographic}{latitudes}. L'OLR, qui varie peu, ne
compense pas la chute de la lumière solaire absorbée.
\end{solution}

\begin{solution}{exo:b1:radiation-heat-transport:4}
(a)~$2.82 \times 365.25 = \qty{1030}{mm}$, \qty{1.03}{m} par an. (b)~$1.030 \times
78.1 = \qty{80.4}{W.m^{-2}}$.
(c)~Le $LE$ de la réanalyse est plus grand de \qty{2}{W.m^{-2}}. Sur de nombreuses années
l'atmosphère ne peut continuer d'accumuler de la vapeur d'eau~: toute l'eau évaporée retombe,
donc $LE$ et $LP$ doivent être égaux en moyenne globale. La différence est
un déséquilibre du modèle, comme son déséquilibre radiatif.
\end{solution}

\begin{solution}{exo:b1:radiation-heat-transport:5}
(a)~$\bar Q = S_0(a/r)^2\sin\delta = 1361 \times 0.9683 \times 0.3978 =
\qty{524.2}{W.m^{-2}}$. (b)~$1361 \times 1.0334 \times 0.3978 =
\qty{559.5}{W.m^{-2}}$~; le rapport est $1.0334/0.9683 = 1.067$. (c)~Le pôle a
de la lumière 24 heures sur 24 à une hauteur constante, donc son $\sin h_\odot$ moyen sur
le jour est $\sin\delta = 0.398$. À l'équateur le Soleil est haut à midi, mais couché
la moitié du temps~; la moyenne de $\sin h_\odot$ sur les 24 heures est $\cos\delta/\pi =
0.292$. Le rapport est $\pi\tan\delta = 1.362$
(\cref{prop:b1:radiation-heat-transport:polar-summer}).
\end{solution}

\begin{solution}{exo:b1:radiation-heat-transport:6}
(a)~$\cos h_0 = -\tan\ang{40}\tan\ang{23.44} = -0.3638$, $h_0 = \ang{111.33} =
\qty{1.943}{rad}$~; le jour dure $24 \times 111.33/180 = \qty{14.84}{h}$~; $\bar Q =
(1361/\pi) \times 0.9683 \times (1.943 \times 0.6428 \times 0.3978 + 0.7660 \times
0.9175 \times 0.9315) = \qty{483.0}{W.m^{-2}}$. (b)~$\cos h_0 = +0.3638$, $h_0 =
\ang{68.67}$, \qty{9.16}{h}~; $\bar Q = \qty{155.9}{W.m^{-2}}$. (c)~Le rapport est
$483.0/155.9 = 3.10$. La distance contribue $0.9683/1.0334 = 0.937$~: elle joue
contre l'été boréal. La géométrie, la hauteur du Soleil et la durée
du jour, contribue $3.10/0.937 = 3.31$, tout le contraste saisonnier et plus.
\end{solution}

\begin{solution}{exo:b1:radiation-heat-transport:7}
(a)~$B = H/LE$~: \num{0.093} à l'équateur, \num{0.374} à \ang{25}~N,
\num{0.338} à \ang{45}~N, \num{-0.152} à \ang{70}~N. (b)~\ang{25}~N~: la moyenne
zonale comprend les grands déserts chauds, terres sèches avec peu d'eau à évaporer, où
le sol chauffé réchauffe l'air à la place. (c)~$H < 0$ signifie que la chaleur passe de
l'air à la surface~: en moyenne annuelle la neige et la glace de mer sont plus froides que l'air
au-dessus d'elles, qui est réchauffé par de l'air venu de \omterm{def:b1:spherical-astronomy:geographic}{latitudes} plus basses.
\end{solution}

\begin{solution}{exo:b1:radiation-heat-transport:8}
(a)~Évalué~: $\alpha = 100/340 = 0.294$, $R = 340 - 100 - 239 =
\qty{+1}{W.m^{-2}}$ (les \num{0.7} arrondis). Réanalyse~: $\alpha = 117.0/341.9 =
0.342$, $R = 341.9 - 117.0 - 237.7 = \qty{-12.8}{W.m^{-2}}$ avec ces termes arrondis,
\qty{-12.9}{W.m^{-2}} avec ceux, non arrondis, de
l'\cref{ex:b1:radiation-heat-transport:global}. (b)~La réanalyse
absorbe $224.8$ contre $240$, \qty{15.2}{W.m^{-2}} de moins, parce qu'elle réfléchit
\qty{17}{W.m^{-2}} de plus~; son OLR n'est inférieur que de \qty{1.3}{W.m^{-2}}. L'excès
de réflexion explique le déséquilibre. (c)~L'intégration de $R$ du pôle Sud au pôle
Nord donne l'intégrale globale, $-12.9 \times 4\pi a_\oplus^2 = -12.9 \times
5.10\times10^{14} = \qty{-6.6}{PW}$. Un transport à travers le pôle Nord est
impossible~: il n'y a rien au-delà, le parallèle est un point. La proposition
ne vaut que si $\int R\,\mathrm dA = 0$, ce qu'exige le régime permanent.
\end{solution}

\begin{solution}{exo:b1:radiation-heat-transport:9}
(a)~L'aire est $2\pi a_\oplus^2 \times 2\sin\ang{30} = 2\pi a_\oplus^2 =
\qty{2.55e14}{m^2}$, la moitié de la Terre~; l'excédent est $33 \times 2.55\times10^{14} =
\qty{8.4}{PW}$. (b)~\qty{4.2}{PW} à travers chaque parallèle. (c)~Les deux font
\qty{8.5}{PW}, l'excédent, mais elles sont inégales~: l'hémisphère sud exporte
\qty{0.5}{PW} de plus que le nord. L'excédent n'est pas partagé à
l'équateur~: dans cette réanalyse \qty{0.36}{PW} de l'excédent des tropiques nord traversent
l'équateur vers le sud.
\end{solution}

\begin{solution}{exo:b1:radiation-heat-transport:10}
(a)~Écrivons la moyenne journalière $\bar Q = (a/r)^2 F(\varphi, \delta)$, avec $F$ la
parenthèse du \cref{thm:b1:radiation-heat-transport:daily} multipliée par $S_0/\pi$. Pendant
$\mathrm dt$ l'énergie reçue par mètre carré est $\bar Q\,\mathrm dt$, et
$\mathrm dt = r^2\,\mathrm d\lambda/h$, donc $\bar Q\,\mathrm dt = (a^2/h)
F(\varphi, \delta(\lambda))\,\mathrm d\lambda$~: le facteur $r^2$ se simplifie.
(b)~Sur une année l'énergie est $(a^2/h)\int_0^{2\pi} F(\varphi,
\delta(\lambda))\,\mathrm d\lambda$. La formule est inchangée quand $\varphi$ et
$\delta$ changent de signe ensemble, $F(-\varphi, \delta) = F(\varphi, -\delta)$, et
$\delta(\lambda + \pi) = -\delta(\lambda)$, puisque $\sin\delta = \sin\varepsilon
\sin\lambda$. Le décalage de $\lambda$ de $\pi$ dans l'intégrale montre que $-\varphi$
reçoit la même énergie que $\varphi$. (c)~Le printemps et l'été boréaux durent
\num{7.6} jours de plus, parce que la Terre est près de l'aphélie et se déplace alors lentement~;
l'été austral est plus lumineux mais plus court. D'après (a), les deux effets se compensent exactement.
\end{solution}

\begin{solution}{exo:b1:radiation-heat-transport:11}
(a)~$\langle (a/r)^2\rangle = \frac1P\int_0^P \frac{a^2}{r^2}\,\mathrm dt =
\frac1P\int_0^{2\pi}\frac{a^2}{r^2}\frac{r^2}{h}\,\mathrm d\lambda = \frac{2\pi
a^2}{hP} = \frac{2\pi a^2}{2\pi a b} = \frac ab$, et $b = a\sqrt{1-e^2}$.
(b)~La moyenne globale instantanée $\frac14 S_0(a/r)^2$ a donc pour moyenne $\frac14 S_0
\langle(a/r)^2\rangle = S_0/(4\sqrt{1-e^2})$. (c)~Numériquement, $340.25/\sqrt{1 - 0.0167^2} =
\qty{340.30}{W.m^{-2}}$, soit \qty{0.05}{W.m^{-2}} ou \qty{0.014}{\percent} au-dessus de
$S_0/4$.
\end{solution}

\begin{solution}{exo:b1:radiation-heat-transport:12}
(a)~Atmosphère $0.78 \times 4.10 = \qty{3.20}{PW}$, océan $0.22 \times 4.10 =
\qty{0.90}{PW}$. (b)~On trouve que l'atmosphère seule porte plus, \qty{5.0}{PW},
que l'atmosphère et l'océan de la réanalyse ensemble, \qty{4.02}{PW}~: le
transport total de la réanalyse est trop faible d'au moins un cinquième ($1 - 4.02/5.0 = 0.20$).
Son \omterm{def:b1:radiation-heat-transport:net-radiation}{rayonnement net} varie trop peu entre les tropiques et les pôles. (c)~Retirer une
constante de $R$ rend son intégrale globale nulle mais laisse inchangée sa variation avec la
\omterm{def:b1:spherical-astronomy:geographic}{latitude}~; l'erreur réside dans cette variation, dans l'albédo des nuages et des surfaces du
modèle à chaque \omterm{def:b1:spherical-astronomy:geographic}{latitude}, et pas seulement dans la moyenne globale.
\end{solution}

\begin{solution}{pb:b1:radiation-heat-transport:1}
\begin{enumerate}
  \item Juin~: $\cos h_0 = -0.3638$, $h_0 = \ang{111.33}$, $\bar Q =
        \qty{483.0}{W.m^{-2}}$. Décembre~: $h_0 = \ang{68.67}$, $\bar Q =
        \qty{155.9}{W.m^{-2}}$. Mars~: $h_0 = \ang{90}$, $\bar Q = (1361/\pi)
        \times 1.0081 \times \cos\ang{40} = \qty{334.6}{W.m^{-2}}$.
  \item En juin~: $\tan\ang{80}\tan\ang{23.44} = 2.46 > 1$, jour polaire, $h_0 = \pi$~:
        $\bar Q = 1361 \times 0.9683 \times \sin\ang{80}\sin\ang{23.44} =
        \qty{516.3}{W.m^{-2}}$. Décembre~: nuit polaire, $\bar Q = 0$. Mars~:
        $(1361/\pi) \times 1.0081 \times \cos\ang{80} = \qty{75.8}{W.m^{-2}}$.
  \item À \ang{80}~N le Soleil ne se couche pas ($h_0 = \ang{180}$, 24 heures) alors qu'à
        \ang{40}~N il brille \qty{14.8}{h}. Sa hauteur est plus basse, entre
        \ang{13.4} et \ang{33.4}, mais le terme $h_0\sin\varphi\sin\delta$, qui
        croît avec la durée du jour et la \omterm{def:b1:spherical-astronomy:geographic}{latitude}, l'emporte.
  \item Les valeurs de bande dépassent les valeurs centrales de \qty{0.5}{\percent},
        \qty{0.8}{\percent}, \qty{1.7}{\percent}, \qty{1.7}{\percent} et
        \qty{1.3}{\percent}. D'abord, la réanalyse utilise un éclairement solaire légèrement plus
        grand (question~5), \qty{0.5}{\percent} de plus. Ensuite, une moyenne de bande
        n'est pas la valeur au centre~: la moitié équatoriale de chaque bande a plus
        d'aire et plus de lumière que la moitié polaire.
  \item D'après la \cref{prop:b1:radiation-heat-transport:global}, $S_0 = 4\sqrt{1-e^2}
        \times 341.9 = \qty{1367}{W.m^{-2}}$, \qty{0.47}{\percent} de plus que \qty{1361}{W.m^{-2}}.
  \item $A = 2\pi a_\oplus^2(1 - \sin\ang{40}) = 2\pi \times (6.371\times10^6)^2
        \times 0.3572 = \qty{9.110e13}{m^2}$, une fraction $(1 - \sin\ang{40})/2 =
        0.179$ de la Terre.
  \item $(1 - \alpha)\bar Q$~: \num{198.0}, \num{159.3}, \num{124.7}, \num{92.5},
        \num{74.7} \unit{W.m^{-2}}.
  \item $R$~: \num{-28.1}, \num{-51.0}, \num{-76.6}, \num{-104.9}, \num{-116.9}
        \unit{W.m^{-2}}. Les cinq bandes sont en déficit~; la bande polaire
        \ang{80}--\ang{90} a le plus grand.
  \item $T_e = (\text{OLR}/\sigma)^{1/4}$~: \num{251.3}, \num{246.8}, \num{244.1},
        \num{242.9}, \num{241.1} \unit{K}.
  \item Absorbé~: $198.0/74.7 = 2.65$~; OLR~: $226.1/191.6 = 1.18$.
  \item Le Soleil est plus bas, et la lumière arrivant sous un angle rasant est davantage réfléchie
        par les nuages et l'eau~; la neige et la glace de mer, surfaces brillantes, couvrent davantage le
        sol et la mer vers le pôle.
  \item L'OLR varie comme $T_e^4$, et $T_e$ ne varie que de \qty{10}{K}, \qty{4}{\percent},
        sur la calotte~: les hautes \omterm{def:b1:spherical-astronomy:geographic}{latitudes} sont maintenues bien plus chaudes que leur lumière solaire
        seule ne le permettrait, par la chaleur apportée du sud.
  \item En \unit{m^2}~: \num{3.143e13}, \num{2.550e13}, \num{1.879e13},
        \num{1.151e13}, \num{0.387e13}~; la somme est \qty{9.110e13}{m^2}.
  \item $R_iA_i$~: \num{-0.883}, \num{-1.300}, \num{-1.439}, \num{-1.207},
        \num{-0.453} \unit{PW}~; la calotte~: \qty{-5.28}{PW}.
  \item $\alpha_\text{calotte} = \sum\alpha_i\bar Q_iA_i/\sum\bar Q_iA_i = 0.400$,
        contre \num{0.342} pour le globe.
  \item $-12.90 \times 5.10\times10^{14} = \qty{-6.58}{PW}$~: la Terre perdrait
        de la chaleur à ce rythme et se refroidirait. Le déséquilibre évalué est
        $+0.7 \times 5.10\times10^{14} = \qty{+0.36}{PW}$, dix-huit fois plus petit et
        de signe opposé~: le déséquilibre de la réanalyse est une erreur de modèle.
  \item $R_i + 12.90$ corrigé~: \num{-15.2}, \num{-38.1}, \num{-63.7},
        \num{-92.0}, \num{-104.0} \unit{W.m^{-2}}~; puissances \num{-0.478},
        \num{-0.972}, \num{-1.197}, \num{-1.059}, \num{-0.403} \unit{PW}~; la calotte~:
        \qty{-4.108}{PW}, un déficit moyen de \qty{45.1}{W.m^{-2}}.
  \item L'énergie de la calotte ne change pas d'une année à l'autre, donc ce qu'elle perd au
        sommet de l'atmosphère doit y entrer, et la seule voie d'entrée est à travers le
        40\textsuperscript{e} parallèle~: $T(\ang{40}) = -2\pi a_\oplus^2\int_{\ang{40}}^{\ang{90}}
        R\cos\varphi\,\mathrm d\varphi$.
  \item $T(\ang{40}\,\text{N}) = +\qty{4.11}{PW}$, vers le nord.
  \item $LP = 2.4655\times10^6 \times 1.95/86400 = \qty{55.6}{W.m^{-2}}$, plus que
        $LE = \qty{42.4}{W.m^{-2}}$~: il tombe plus d'eau sur la calotte qu'il n'en
        évapore~; la différence arrive sous forme de vapeur.
  \item $(LE - LP) - 2.07 = 42.4 - 55.6 - 2.07 = \qty{-15.3}{W.m^{-2}}$~; multiplié par
        l'aire de la calotte~: \qty{-1.40}{PW}. L'atmosphère apporte \qty{1.40}{PW}
        de chaleur latente à travers \ang{40}~N.
  \item $1.40/4.11 = 0.34$~: un tiers.
  \item Océan $0.22 \times 4.11 = \qty{0.90}{PW}$~; atmosphère $0.78 \times 4.11 =
        \qty{3.20}{PW}$, dont \qty{1.40}{PW} de chaleur latente et \qty{1.80}{PW}
        portés sous forme de chaleur sensible et d'énergie potentielle de l'air.
  \item Une borne inférieure. Le total de la réanalyse à \ang{43}~N, \qty{4.02}{PW}, est inférieur à
        l'atmosphère seule dans l'étude, \qty{5.0}{PW}~; le vrai total est
        plus grand.
  \item Environ \qty{4.11}{PW}, \qty{4.11e15}{W}, d'après le rayonnement de la
        réanalyse, dont un tiers sous forme de chaleur latente~; la vraie valeur est plus grande encore.
\end{enumerate}
\end{solution}
